\documentclass[11pt]{article}

\usepackage[utf8]{inputenc}
\usepackage[T1]{fontenc}
\usepackage{lmodern}
\usepackage{geometry}
\usepackage{amsmath,amssymb,amsfonts,amsthm,mathtools}
\usepackage{bm}
\usepackage{enumitem}
\usepackage{microtype}
\usepackage{hyperref}
\usepackage{titlesec}
\usepackage{booktabs}
\usepackage{array}
\usepackage{setspace}
\usepackage{mathrsfs}
\usepackage{physics}

\geometry{margin=1in}
\hypersetup{colorlinks=true, linkcolor=black, urlcolor=blue, citecolor=black}
\setlist[enumerate]{topsep=0.4em,itemsep=0.35em}
\setlist[itemize]{topsep=0.3em,itemsep=0.25em}
\titleformat{\section}{\large\bfseries}{\thesection.}{0.5em}{}
\titleformat{\subsection}{\normalsize\bfseries}{\thesubsection.}{0.5em}{}
\setstretch{1.06}

\newcommand{\Aether}{\AE{}ther}
\newcommand{\Aflow}{\AE{}ther-flow}
\newcommand{\TheoryName}{\Aflow{} Interpretation of Relativity}
\newcommand{\TheoryProgram}{\Aether{} / \Aflow{} framework}
\newcommand{\CompletedMetric}{g^{\sharp}}
\newcommand{\CompletionCore}{\mathfrak C^{\mathrm{zr,aff}}}
\newcommand{\TraceCore}{\mathfrak C^{\mathrm{tr}}}
\newcommand{\AffineNucleus}{\mathfrak N^{\mathrm{aff}}}
\newcommand{\AmbientPackage}{\mathfrak A^{\mathrm{bal}}}
\newcommand{\CombinedLaw}{\mathfrak R^{\mathrm{comb}}}
\newcommand{\SelectorLaw}{\mathfrak S^{\mathrm{sel},0}}
\newcommand{\SelectorCore}{\mathfrak C^{\mathrm{sel,split}}}
\newcommand{\ShellRestriction}{\mathfrak S^{\mathrm{sub,sh}}}
\newcommand{\LiftPackage}{\mathfrak L^{\mathrm{sub,lift}}}
\newcommand{\PrecursorPackage}{\mathfrak P^{\mathrm{sub,prec}}}
\newcommand{\OriginPackage}{\mathfrak O^{\mathrm{sub,orig}}}

\newtheorem{definition}{Definition}
\newtheorem{proposition}{Proposition}
\newtheorem{remark}{Remark}
\newtheorem{corollary}{Corollary}
\newtheorem{theorem}{Theorem}

\title{\textbf{The \TheoryName{}}\\[0.4em]
\large Primitive-Reservoir Observer-Reduction\\
\large Lift-Generating Substrate Precursor Dynamics\\
\large Next Benchmark Criterion}
\author{}
\date{}

\begin{document}

\maketitle

\begin{abstract}
The current lift-frontier routing note had fixed the admissible burden order below the theorem forcing the shell-generating substrate lift package \(\LiftPackage\) from lift-generating substrate precursor dynamics \(\PrecursorPackage\): first try to derive or uniquely force the precursor package \(\PrecursorPackage\) from still more primitive explicit substrate structure, and if that could not yet be done cleanly, prove a genuinely smaller theorem-level collapse or sharper necessity / uniqueness result strictly inside lift dynamics. The new precursor-forcing theorem has now spent that first admissible route. The lift-generating substrate precursor package \(\PrecursorPackage\) is no longer left as an independently inserted stopping point; it is a canonical and unique output of precursor-generating substrate origin dynamics \(\OriginPackage\).

This note records the routing consequence of that new theorem. The next honest benchmark-facing manuscript must therefore go one level below the current precursor frontier: first derive or uniquely force the precursor-generating substrate origin package \(\OriginPackage\) from still more primitive explicit substrate structure in a way that actually changes benchmark-facing status, or, if that cannot be done cleanly, prove a genuinely smaller theorem-level collapse, sharper necessity result, or sharper uniqueness result strictly inside the lift-generating substrate precursor package \(\PrecursorPackage\). Another shell-restriction relay, split-core relay, selector-law restatement, combined-response restatement, ambient-package restatement, trace-core restatement, completion-core restatement, exact-shell affine nucleus restatement, another lift relay, or another precursor relay that preserves the same \(\PrecursorPackage\) without shrinking or forcing it further does not count next.

The benchmark boundary remains fixed throughout. The one operative metric is still exactly \(\CompletedMetric\), the ordinary matter route is unchanged, there is no second operative metric, the low-energy matter law is not enlarged, and stronger first-principles promotion language remains paused. The positive result of this note is therefore routing discipline at the new precursor frontier.
\end{abstract}

\section{Introduction}

The active derivational program of \TheoryName{} is no longer at a stage where the lift-generating substrate precursor package \(\PrecursorPackage\) should be treated as the default unexplained stopping point. The lift-forcing theorem had already shown that the shell-generating substrate lift package \(\LiftPackage\) is a canonical and unique output of \(\PrecursorPackage\). The new precursor-forcing theorem has now spent the next admissible upstream move below that frontier: \(\PrecursorPackage\) itself is a canonical and unique output of the deeper origin class \(\OriginPackage\).

The present note records the routing consequence of that theorem. It does not reopen the precursor mathematics, and it does not claim a completed first-principles substrate derivation of gravity. It records only which kind of next manuscript would now count as an honest benchmark-facing gain below the current precursor frontier.

\paragraph{Framework and claim boundary.}
\TheoryProgram{} denotes the broader \Aether{} / \Aflow{} framework built on the claims that the \Aether{} is the underlying four-dimensional substrate of reality, the \Aflow{} is its intrinsic ordered motion, observed three-dimensional space is the local experiential slice of that deeper substrate, S-time is the experienced order of change arising from matter, light, and the \Aflow{}, and observed expansion is the three-dimensional appearance of deeper four-dimensional ordered motion. Within that framework, \TheoryName{} denotes the exact-closure theory in which the effective gravitational dynamics are adopted to be exactly Einsteinian with universal matter coupling. In the active sequence, \emph{adoption} means use of that established relativistic dynamics without claiming substrate derivation, while \emph{derivation} is reserved for a first-principles recovery from explicit substrate variables.

The present note stays strictly inside that boundary. It records only which kind of next manuscript would now count as an honest benchmark-facing gain below the precursor frontier.

\section{Fixed Inputs}

\begin{definition}[Precursor-frontier fixed inputs]
The criterion recorded here uses the following fixed inputs.
\begin{enumerate}
    \item The benchmark exact-closure package, which fixes the public one-metric GR target of the repository.
    \item The benchmark gatekeeping layer, including the post-bridge status correction and the upstream-compression safeguard that blocks benchmark-neutral deeper-origin relays from counting as new front-line progress once the benchmark-relevant object is unchanged.
    \item The preserved lift-forcing theorem, which already shows that the shell-generating substrate lift package \(\LiftPackage\) is a canonical and unique output of lift-generating substrate precursor dynamics \(\PrecursorPackage\).
    \item The preserved precursor-origin theorem, which already shows that the lift-generating substrate precursor package \(\PrecursorPackage\) can be recovered canonically from precursor-generating substrate origin dynamics \(\OriginPackage\).
    \item The new precursor necessity / uniqueness theorem, which proves that \(\PrecursorPackage\) itself is a canonical and unique output of \(\OriginPackage\).
\end{enumerate}
\end{definition}

\begin{remark}
The present note does not replace those manuscripts. It records the routing consequence of reading them together under the active safeguard against recursive depth drift.
\end{remark}

\section{Why the Precursor Theorem Is the Frontier}

\begin{definition}[Benchmark-neutral precursor relay]
A \emph{benchmark-neutral precursor relay} is a manuscript that preserves the same benchmark one-metric output, the same ordinary matter route, the same lift-generating precursor package \(\PrecursorPackage\), the same shell-generating lift package \(\LiftPackage\), the same shell-restriction package \(\ShellRestriction\), the same split selector-shell core \(\SelectorCore\), the same selector-law consequence \(\SelectorLaw\), the same combined-response consequence \(\CombinedLaw\), the same ambient-package consequence \(\AmbientPackage\), the same trace-core consequence \(\TraceCore\), the same completion-core consequence \(\CompletionCore\), the same exact-shell affine nucleus consequence \(\AffineNucleus\), and the same claim boundary while merely restating \(\PrecursorPackage\), \(\LiftPackage\), \(\ShellRestriction\), \(\SelectorCore\), \(\SelectorLaw\), \(\CombinedLaw\), \(\AmbientPackage\), \(\TraceCore\), \(\CompletionCore\), \(\AffineNucleus\), or the same origin consequence without shrinking or forcing the precursor package further.
\end{definition}

\begin{proposition}[The precursor-forcing theorem is the live frontier]
At the current stage of the primitive-reservoir observer-reduction line, the theorem forcing the lift-generating substrate precursor package \(\PrecursorPackage\) from precursor-generating substrate origin dynamics \(\OriginPackage\) is the live benchmark-facing frontier.
\end{proposition}

\begin{proof}
That theorem changes benchmark-facing status at current scope because it removes the precursor package \(\PrecursorPackage\) itself as the live internal stopping point. The current bridge datum is no longer merely a deeper package whose lift consequence was already known; it is now forced from the origin layer. By contrast, a manuscript that preserves the same benchmark package, the same precursor consequence, and the same lift/shell-restriction/split-core consequences while merely restating the precursor package or another already-forced larger or smaller bridge object does not alter benchmark-facing status unless it adds a comparably sharp gain. Therefore the live frontier is the precursor-forcing theorem rather than the earlier lift-forcing theorem or a later same precursor relay.
\end{proof}

\begin{corollary}[Status of the precursor package as a next-step target]
Under the current routing safeguard, the lift-generating substrate precursor package \(\PrecursorPackage\) is no longer itself the default next benchmark-facing target.
\end{corollary}

\begin{proof}
The precursor-forcing theorem already proves that \(\PrecursorPackage\) is a canonical and unique output of \(\OriginPackage\). Once that is on record, another manuscript that spends the move on the precursor package again does not remove any further benchmark-relevant freedom. The open burden has therefore moved below precursor scope.
\end{proof}

\section{The Next Benchmark Criterion}

\begin{definition}[Admissible next benchmark-facing precursor manuscript]
At the current precursor frontier, a manuscript counts as an admissible next benchmark-facing move only if it respects the following burden order:
\begin{enumerate}
    \item its primary attempt is to derive or uniquely force the precursor-generating substrate origin package \(\OriginPackage\) from still more primitive explicit substrate structure in a way that does more than merely relocate the unexplained layer deeper;
    \item if that cannot be done cleanly at current scope, its admissible fallback gain is to prove a genuinely smaller theorem-level collapse, sharper necessity result, or sharper uniqueness result strictly inside the lift-generating substrate precursor package \(\PrecursorPackage\);
    \item in either case it preserves the same benchmark one-metric package, the same ordinary matter route, and the same claim boundary, and it introduces neither a second operative metric nor an enlarged low-energy matter law and licenses no stronger first-principles promotion language.
\end{enumerate}
\end{definition}

\begin{theorem}[Next honest benchmark-facing task at precursor scope]
The next honest benchmark-facing manuscript at the current precursor frontier should therefore go one level below that frontier: first try to derive or uniquely force the precursor-generating substrate origin package \(\OriginPackage\) from still more primitive explicit substrate structure in a way that actually changes benchmark-facing status. If that cannot be done cleanly, the admissible fallback is to prove a genuinely smaller collapse, sharper necessity result, or sharper uniqueness result strictly inside the lift-generating substrate precursor package \(\PrecursorPackage\). A manuscript that only restates the shell-restriction package \(\ShellRestriction\), only restates the split selector-shell core \(\SelectorCore\), only restates the selector law \(\SelectorLaw\), only restates the combined response law \(\CombinedLaw\), only restates the ambient package \(\AmbientPackage\), only restates the trace core \(\TraceCore\), only restates the completion-core presentation \(\CompletionCore\), only restates the exact-shell affine nucleus \(\AffineNucleus\), only restates the shell-generating substrate lift package \(\LiftPackage\), or only repeats the same precursor consequence while preserving \(\PrecursorPackage\) without shrinking or forcing it further does not satisfy the criterion.
\end{theorem}

\begin{proof}
By Proposition 1, the live frontier already lies at the theorem that forces the precursor package from origin dynamics. Therefore a second manuscript below it must improve on that status rather than merely accompany it. A deeper-origin manuscript can do so only if it changes benchmark-facing status by more than depth alone, and the first clean way to do that is to force \(\OriginPackage\) itself from still more primitive explicit substrate structure. If that first option is not yet available cleanly, the routing safeguard allows the fallback of a smaller internal collapse, necessity result, or uniqueness result inside \(\PrecursorPackage\). If a manuscript simply reproduces the same precursor consequence through another relay while preserving the same precursor package without shrinking or forcing it further, the routing rule classifies it as benchmark neutral.
\end{proof}

\section{What the Next Move Should Not Be}

\begin{proposition}[Screened next moves]
At the current precursor frontier, none of the following is the default next benchmark-facing move:
\begin{enumerate}
    \item another restatement of the primitive substrate shell-restriction package \(\ShellRestriction\);
    \item another restatement of the split zero-hidden selector-shell core \(\SelectorCore\);
    \item another restatement of the full zero-hidden selector law \(\SelectorLaw\);
    \item another restatement of the combined response substrate law \(\CombinedLaw\);
    \item another restatement of the completion-balance ambient dynamics package \(\AmbientPackage\);
    \item another restatement of the completion-state shell-trace core \(\TraceCore\);
    \item another restatement of the full completion-core presentation \(\CompletionCore\);
    \item another restatement of the exact-shell affine nucleus \(\AffineNucleus\);
    \item another restatement of the shell-generating substrate lift package \(\LiftPackage\);
    \item another precursor relay that preserves the same \(\PrecursorPackage\) without shrinking or forcing it further;
    \item a manuscript that introduces a second operative metric, enlarges the ordinary matter law, or uses stronger first-principles promotion language.
\end{enumerate}
\end{proposition}

\begin{proof}
Items (1)--(10) fail the preceding criterion because they do not directly address the current origin-facing burden. They revisit either already-forced larger bridge objects, already-forced smaller bridge objects, or the same precursor stopping point rather than removing the remaining precursor freedom or proving a still smaller internal datum inside \(\PrecursorPackage\). Item (11) violates the fixed claim boundary of the benchmark package and therefore cannot count as a disciplined next step on the active line. The benchmark package remains the exact one-metric GR package already fixed by adoption, so any admissible next manuscript must preserve that boundary while keeping stronger first-principles promotion language paused.
\end{proof}

\begin{remark}
This screening statement is not a denial that the shell-restriction forcing theorem, the lift-forcing theorem, the selector-law collapse theorem, the selector-law forcing theorem, the combined-response theorem, the ambient-package theorem, the trace-core theorem, or the completion-core theorem matter historically. It is a routing statement: they are not the honest way to spend the next benchmark-facing move from the current precursor frontier.
\end{remark}

\section{Conclusion}

The positive result of this note is routing discipline at the current precursor frontier. The precursor-forcing theorem remains the live benchmark-facing gain because it already removes the lift-generating substrate precursor package as the internal stopping point on the active line. The precursor package therefore no longer sets the honest next burden, and neither do the already-forced lift package, shell-restriction package, split core, selector law, combined response law, ambient package, trace core, completion-core presentation, or exact-shell affine nucleus.

The scope remains narrow and honest. The benchmark package is unchanged: the one operative metric is still exactly \(\CompletedMetric\), the ordinary matter route is unchanged, there is no second operative metric, and the low-energy matter law is not enlarged. Nothing here upgrades adoption to derivation or licenses stronger first-principles promotion language yet.

The next open burden is therefore clear. The next benchmark-facing manuscript should go one level below the current precursor frontier: first try to derive or uniquely force the precursor-generating substrate origin package \(\OriginPackage\) from still more primitive explicit substrate structure in a way that actually changes benchmark-facing status while preserving the same benchmark one-metric package and claim boundary. If that cannot be done cleanly, the admissible fallback is a genuinely smaller collapse, sharper necessity result, or sharper uniqueness result strictly inside the lift-generating substrate precursor package \(\PrecursorPackage\). Another shell-restriction relay, split-core relay, selector-law restatement, combined-response restatement, ambient-package restatement, trace-core restatement, completion-core restatement, exact-shell affine nucleus restatement, another lift relay, or another precursor relay that preserves the same precursor package without shrinking or forcing it further is not the clean next manuscript target at current scope.

\end{document}
